\section{\ref{sec:dishbrain}장. DishBrain}

테스트 벤치의 구조와 게임 결과는 실험 논문 하나와 그 후속 연구에서 가져왔다. 잡음 공식과 좋은 설정 쪽으로의 표류는 본문에서 유도하고 이 책의 모델로 확인했으며, 접시 안의 학습 메커니즘은 밝혀지지 않았다.

{\footnotesize
\setlength{\tabcolsep}{3pt}\renewcommand{\arraystretch}{1.15}
\begin{longtable}{>{\raggedright}p{0.5cm}>{\raggedright}p{4.0cm}>{\raggedright}p{5.6cm}>{\raggedright}p{2.3cm}>{\raggedright\arraybackslash}p{1.7cm}}
\toprule
번호 & 주장 & 실험과 측정된 것 & 출처 & 상태 \\
\midrule\endfirsthead
\toprule
번호 & 주장 & 실험과 측정된 것 & 출처 & 상태 \\
\midrule\endhead
1 & 테스트 벤치: 칩당 생쥐 피질 세포 약 $800\,000$개 또는 인간 세포 약 $10^6$개. $8$~mm$^2$에 전극 $26\,000$개, 기록 최대 $1024$, 자극은 $8$개로 & 논문의 방법: MaxOne 칩, $8$~mm$^2$에 백금 전극 $26\,000$개, 초당 $20\,000$ 표본으로 $1024$채널 기록. 자극은 기술적으로 $32$개 전극까지 가능하나 독립적으로는 $8$개. 생쥐 배양은 약 $10$일째부터 사용 & \cite{Kagan2022} & 측정됨 \\
2 & $10$~ms 클록의 루프: 운동 영역의 스파이크가 라켓을 움직인다(그림~\ref{fig:dishloop}) & 스파이크를 $10$~ms 창으로 센다. 스파이크에서 응답 자극까지의 지연은 약 $5$~ms로, 장비의 버퍼가 정한다 & \cite{Kagan2022} & 측정됨 \\
3 & 입력: 장소 부호($8$개 전극 중 어느 것)와 초당 $4$--$40$회의 빈도 부호 & 본 실험에서 자극 빈도는 공이 먼 벽에 있을 때 $4$~Hz에서 라켓 앞 $40$~Hz까지 선형으로 오른다. 공 펄스의 진폭은 $75$~mV, 펄스는 직사각형 2상 & \cite{Kagan2022} & 측정됨 \\
4 & 출력 $\Delta p=c\,(n_\uparrow-n_\downarrow)$~\eqref{eq:dishreadout}, 창당 스파이크 수의 잡음 $1/\sqrt{RT}$ & 두 영역의 차이 규칙은 논문에 기술되어 있으나(활동에 따른 영역 가중치 포함) 정확한 공식은 없다. 잡음 추정은 푸아송 계수의 결과이며 본문에서 유도했다 & \cite{Kagan2022}; 본문의 유도 & 이론 \\
5 & 교사: 히트 뒤에는 예측 가능한 자극 $75$~mV, 초당 $100$회, $0.1$~s. 미스 뒤에는 예측 불가능한 자극 $150$~mV, 초당 $5$회, $4$~s 동안 무작위 장소와 시점에, 그 뒤 $4$~s 휴지 & 방법: 히트 뒤 $75$~mV, $100$~Hz, $100$~ms를 $8$개 전극 모두에 동시에. 미스 뒤 $150$~mV, $5$~Hz를 $4$~s 동안 무작위 전극에 무작위 시점으로, 그 뒤 $4$~s 자극 없음. 배양에 도파민은 없다 & \cite{Kagan2022} & 측정됨 \\
6 & 네트워크는 자기 입력이 예측 가능해지도록 행동한다 & 자유 에너지 원리는 저자들이 프로토콜을 이끌어 낸 이론적 제안이다. 실험은 이와 맞지만 더 단순한 메커니즘(7--8행)과 구별하지 못한다 & \cite{Friston2010,Kagan2022} & 가설 \\
7 & 실패가 네트워크를 흔들고 성공은 흔들지 않으면, 설정에 머무는 시간은 $\rho\propto1/P_{\text{miss}}$~\eqref{eq:dishdensity}(명\ref{prop:dishscramble}) & 대칭적 흔들기를 가진 마르코프 연쇄의 흐름 균형에서 나오며, 본문에서 유도했다. 배양에서 직접 확인한 것은 없다 & 본문의 유도 & 이론 \\
8 & `미스 때 뒤섞기' 모델이 학습한다: 히트 비율 $0.21\to0.38$. 모든 랠리 뒤 흔들기는 $\approx0.12$, 흔들기 없음은 $0.19$(그림~\ref{fig:dishmodel}) & 계산, 스크립트 \texttt{models/dishbrain\_model.py}, 모델 배양 $40$개에 랠리 $2000$번씩: $0.21\to0.38$; $0.13\to0.11$; $0.19\to0.19$ & --- & 모델 \\
9 & 랠리가 길어지는 것은 완전한 루프의 살아 있는 배양뿐이고, 대조군은 학습하지 않는다 & $20$~분 세션 $399$번: 생쥐 배양 $9$개(세션 $101$번), 인간 $11$개($138$), 배지만 있는 칩 $6$개($80$), 휴지 배양 $20$개($42$), 모의실험 $3$개($38$). 마지막 $15$~분의 평균 랠리 길이가 처음 $5$~분보다 긴 것은 살아 있는 배양뿐이다. 뉴런이 아닌 세포(HEK293T)와 배지만 있는 칩에서는 효과가 없다 & \cite{Kagan2022} & 측정됨 \\
10 & 세션 안의 유의미한 향상은 완전한 피드백일 때만. 자극 대신 침묵을 주면 세션 끝의 게임은 피드백 없음보다 낫지만 효과가 더 작다 & $3$일 동안 세션 $486$번: `자극'($n=27$), `침묵'($17$), `피드백 없음'($15$). 시간에 따른 향상은 `자극' 조건에서만 유의미하다. 세션 끝에 `침묵' 조건은 더 작은 효과로 `피드백 없음'을 앞섰다 & \cite{Kagan2022} & 측정됨 \\
11 & 효과는 작다. 네트워크가 오래 가는 학습을 하는지는 열린 질문 & 세션 안의 향상은 확실하지만, 저자들 스스로 말하듯 며칠에 걸친 세션 간 학습은 안정적으로 관찰되지 않았다 & \cite{Kagan2022} & 측정됨 \\
12 & 게임 중 네트워크는 임계 상태에 다가가고, 가까울수록 게임을 잘한다 & 같은 배양의 활동 사태 분석: 구조화된 입력은 네트워크를 임계 상태 쪽으로 옮기고, 게임 성적은 임계 상태와의 근접도와 상관된다. 그러나 행동의 결과에 대한 정보 없이 임계성만으로는 학습에 부족하다 & \cite{Habibollahi2023} & 측정됨 \\
13 & 배양의 `지각력'은 보이지 않았다 & 연구자 30명의 반론 논문: 닫힌 루프 안의 행동 변화는 지능도 감각도 증명하지 않는다. 그것을 보여 주는 실험은 없다 & \cite{Balci2023} & 가설 \\
14 & 제안하는 네 번째 대조 실험: 미스 뒤에 같은 길이와 같은 에너지의 예측 가능한 자극(\ref{sec:dishspec}절) & 이런 실험은 수행되지 않았다. 이 책의 제안이다 & --- & 가설 \\
\bottomrule
\end{longtable}}
